\section{Conclusiones}
\label{rz:conclusiones}

La forma de Weil admite, sobre cada ventana fijada, una
reducción exacta a un problema finito que conserva su
complemento infinito. El espacio de detalles de media nula
puede hacerse coercivo mediante refinamiento; su eliminación
produce un complemento de Schur con un criterio equivalente de
positividad y el mismo índice negativo que la forma completa. En la ventana
$(-1/18,1/18)$, el control conjunto de medias, detalles y
acoplamientos demuestra
\[
\mathscr W(f,f)\geq\frac12\|f\|_2^2,
\qquad f\in V_{1/9}.
\]
La afirmación comprende el cierre de las pruebas suaves
compactamente soportadas en la norma conjunta de los lectores.
Es un resultado sobre un dominio funcional infinito, estable
bajo traslación y refinamiento interior.

\subsection*{Signo de la forma y conservación del complemento infinito}

La identidad
\[
\mathscr W_R(Ua+d,Ua+d)
=a^*\bigl(A-B^*D^{-1}B\bigr)a
+\mathfrak d[d+D^{-1}Ba,d+D^{-1}Ba]
\]
explica qué conserva la reducción. El segundo sumando es
positivo y contiene todos los detalles; el primero conserva
su interacción con las medias por medio de $B^*D^{-1}B$.
El teorema~\ref{schur-add-schur-exacto} prueba que la
positividad en $V_R$ equivale a la positividad de esa matriz.
Su tamaño finito no reemplaza el cálculo de sus entradas:
la certificación residual de
\eqref{schur-add-eq-S7} proporciona cotas con el error del
complemento retenido.

En la ventana de longitud $1/9$, las cotas
$D>I$, $A>97/100$ y $\|B\|^2<1/5$ producen el margen
coercivo anunciado. Este caso local no contiene una traslación
prima activa porque $1/9<\log2$. La identidad de transporte
de la forma completa permite, sin embargo, llevar la misma
cota a cada subespacio $C_a^n\tau_t C_c^\infty((-1/18,1/18))$,
con $a>b>1/2$. Sus elementos pueden tener colas exponenciales,
y su contribución prima se conserva con la suma convergente
completa. La estimación se aplica a cada colección entera, no
sólo a una prueba seleccionada.

Los otros dominios positivos mantienen su contenido propio:
el ensamblaje especificado de nueve intervalos controla sus
productos cruzados y todos sus refinamientos, mientras una
ventana arbitraria admite un umbral de coercividad para los
detalles celulares. La distinción entre estos dominios permite
identificar tanto el margen demostrado como el acoplamiento
que decide el signo de una composición más amplia.

\subsection*{La aritmética determina los modos y sus lectores}

El antecedente aritmético está construido antes de la lectura
espectral. Las dos hojas de APP proporcionan las operaciones
locales y sus cocientes; el producto nativo propaga esos datos
en palabras de forma normal única. Las aperturas del producto
seleccionan irreducibles, y la evaluación los identifica con
los primos positivos. La realización operatoria del selector
cuenta las descomposiciones no triviales. Su positividad es
una norma cuadrada con ese significado aritmético concreto.

La factorización organiza después las ocupaciones de modos y
sus relojes. El producto de Euler y los momentos
primo--potencia proceden de las trazas de esa representación.
El límite de la traza de calor cíclica y su transformación de
Mellin producen la función completada, con sus factores gamma
y polares y su ecuación funcional. Las partes finitas, las
evaluaciones regularizadas y las cotas de resto conservan
las normalizaciones necesarias para componer esas lecturas.

La unidad del recorrido reside en los estados enriquecidos
de APP--TRIT--TPK y en la estructura conjunta del continuo:
ruta, orientación, hojas, memoria e incidencias se transmiten
a través de restricciones y refinamientos compatibles. La
evaluación real obtiene un resultado de esa construcción;
la historia y la frontera conservan la información que sus
realizaciones posteriores necesitan.

Las leyes de composición precisan esta continuidad. El producto de
dos momentos se realiza sobre pares de modos y sus coeficientes se
agrupan por convolución de divisores; la composición diagonal mantiene
un solo modo. La normalización espectral transporta toda la expansión
de Laurent, incluido el término polar, mediante una transformación
triangular afín invertible. La ecuación funcional relaciona las
derivadas negativas pares con los valores impares positivos y, en
particular, establece $Z(3)=4\pi^2\log A_2$. Finalmente, los cortes
por irreducibles aportan encierros explícitos complementarios a los
cortes por profundidad. Cada identidad conserva la operación que la
produce y el dominio en que se demuestra.

\subsection*{Frontera, vacancias y doble círculo}

La reconstrucción del interior desde la frontera levantada
en la coordenatización multiafín proporciona una instancia explícita de
conservación. La diferencia suma--producto conserva el dato
residual y el cociente; las lecturas $72,27,77,22$ y la
completación $729+271$ mantienen sus operaciones y dominios.
El centro $(5,5)$ y sus vacancias conjugadas conservan suma
y residuo multiplicativo mientras cambia el cociente. La
firma regional $555555$ requiere tres hojas de continuación,
y el calendario de capacidad produce las longitudes $21,22$
y separaciones $6,7$ con un balance exacto.

El doble círculo lleva esa conservación al transporte de
lectores. La compresión de nueve fases produce el complemento
$(8z,-z,\ldots,-z)/27$ y el factor $8/81$. La razón modal
$q=5/9$ determina los momentos de memoria $q^{|n|}$; su
factorización triangular prueba positividad en toda dimensión
y el determinante $(56/81)^N$. El registro de pérdidas
conserva el terminal y los productos cruzados. El refinamiento
puede conducir a un límite débil atómico de estados, cuya
representación se distingue del desplazamiento regular en
el que se construyeron los vectores iniciales.

Para los lectores aritméticos, la cola gamma reconstruye la
prueba y las filas de ambas columnas. El conector resultante
conserva exactamente
$J_+^*J_+-J_-^*J_-$, con los términos primo, gamma y polar.
Su acción permite formular el problema energético sobre un
operador explícito. La invariancia bajo $C_a$ preserva además
todos los momentos cruzados de la forma, no sólo sus
evaluaciones diagonales.

\subsection*{Condición global y relación con la hipótesis de Riemann}

El criterio global de Weil requiere positividad sobre todo
el dominio de prueba. En la formulación de las dos columnas,
la condición exacta es la contractividad de la transferencia
alcanzable sobre la clausura de la imagen positiva, con
compatibilidad al ampliar las ventanas. En la formulación
circular, se expresa por la identidad
\[
m_n^{\rm ar}(f,g)
=\langle U_{\rm mem}^{\,n}\mathcal E f,
          \mathcal E g\rangle_{\mathcal K_{\rm lim}},
\]
para todos los momentos mixtos y una colección que cubra el
dominio mediante los límites declarados. La identidad de
energía y el entrelazamiento deben darse conjuntamente; la
positividad del Gram regular de memoria no los sustituye.

Las pruebas de este artículo establecen la reconstrucción,
la invariancia y la coercividad en los dominios indicados.
La identificación global anterior continúa siendo una
hipótesis de los resultados espectrales que la emplean y no
se deduce de la positividad de cada colección transportada por
separado. Por ello, el artículo no concluye una resolución
de la hipótesis de Riemann. Su contribución funcional es
precisa: transforma la conservación de la estructura
aritmética y de la memoria en operadores recuperables,
reducciones exactas y cotas coercivas, manteniendo visibles
los productos mixtos y los dominios que gobiernan cualquier
prolongación global.
